\documentclass[11pt]{article}
\usepackage[T1]{fontenc}
\usepackage{amsmath,amssymb,geometry,hyperref}
\geometry{margin=1in}
\title{A singular solvable free Pick problem\\Disproof of Conjecture 00000002352}
\author{ziangni-sys}
\date{October 8, 2026}
\begin{document}
\maketitle
\paragraph{Claim and scope.}
The conjecture identifies solvability of free Nevanlinna--Pick interpolation
on the matrix ball with positive definiteness of its free Pick matrix.
The Chinese version also says positive definite (rather than positive
semidefinite). This strict requirement is false. We give distinct nodes
and strictly interior target values, so no boundary target is involved.
We use the standard Schur-class condition of multiplier norm at most one.
No uniform norm strictly less than one is required in the conjecture.

\paragraph{The genuine free interpolant.}
Take the one-variable free matrix ball
\[
 \mathfrak B_1=\coprod_{n\geq1}\{X\in M_n(\mathbb C):\|X\|_{\mathrm{op}}<1\}.
\]
The free polynomial $F(X)=X$ is defined at every matrix size. It is
analytic, respects direct sums and similarities, and even respects every
intertwiner: $SX=YS$ implies $SF(X)=F(Y)S$. Moreover
$\|F(X)\|_{\mathrm{op}}=\|X\|_{\mathrm{op}}<1$ throughout the open ball.
Its multiplier on one-variable Fock space is the unilateral shift,
an isometry, so its multiplier norm is one. Thus it is a legitimate
contractive free Schur interpolant, including pointwise strict contractivity.
If the free ball has a prescribed number $d>1$ of coordinates, use
$F(X)=X_1$ and nodes $(0,\ldots,0)$ and $(1/2,0,\ldots,0)$ instead;
the same calculation applies.

\paragraph{Data and the actual kernel.}
At scalar level choose $z_1=0$, $z_2=1/2$ and prescribe
$w_1=0$, $w_2=1/2$. These are distinct interior nodes and interior
targets, and $F(z_i)=w_i$. The free Szeg\H{o} kernel in one variable is
\[
 K(Z,W)[P]=\sum_{k\geq0}Z^kP(W^*)^k.
\]
At scalar level its coefficient is the convergent geometric series
$K(z,\zeta)=\sum_{k\geq0}(z\overline\zeta)^k
=(1-z\overline\zeta)^{-1}$. Consequently the actual Pick matrix is
\[
 P_{ij}=(1-w_i\overline{w_j})K(z_i,z_j)
       =\frac{1-z_i\overline{z_j}}{1-z_i\overline{z_j}}=1,
 \qquad
 P=\begin{pmatrix}1&1\\1&1\end{pmatrix}.
\]
All denominators are nonzero (they are $1$ or $3/4$).
For the nonzero vector $v=(1,-1)^T$, $Pv=0$, hence $v^*Pv=0$.
The matrix is not positive definite, even though the interpolation is
solvable. It is positive semidefinite, since
$u^*Pu=|u_1+u_2|^2\geq0$.

\paragraph{Conclusion and interpretation.}
This disproves the positive-definite necessity asserted in the first
conjunct; no disproof of the Szeg\H{o}-kernel identification is needed.
Positive semidefiniteness (or complete positivity in the operator-valued
formulation) is the usual non-strict Schur-class criterion. If
``positive definite'' were intended as terminology for nonnegative
kernels, the statement would need that clarification; our disproof
addresses the strict matrix meaning stated in both languages.
Requiring a uniformly strictly contractive multiplier would change the
problem and exclude this extremal interpolant.

\paragraph{Formal verification.}
The Lean project uses complex Euclidean spaces and continuous linear
endomorphisms, with their operator norm, to verify analyticity,
contractivity and nc intertwining at every size. It also verifies the
Szeg\H{o} geometric-series identity, interior distinct scalar nodes,
interpolation, all four Pick entries, and failure of
\texttt{Matrix.PosDef} using the exhibited nonzero vector.
The finite matrix data are an allowed level of the free problem;
the interpolant is defined on every level. No scalar surrogate is
assumed in place of the free function.

\paragraph{Reference.}
J. A. Ball, G. Marx and V. Vinnikov, \emph{Interpolation and
transfer-function realization for the noncommutative Schur--Agler class},
\href{https://arxiv.org/abs/1602.00762}{arXiv:1602.00762}.
This is background for the standard complete-positivity formulation;
the counterexample and kernel calculation above are self-contained.
\end{document}
